%=====================================================================
%  SSC0903 - Computacao de Alto Desempenho
%  TB1 - Simulacao paralela da propagacao direcional de incendio
%
%  Compilar com:  pdflatex main.tex  (duas vezes, por causa das refs)
%=====================================================================
\documentclass[11pt,a4paper]{article}

\usepackage[utf8]{inputenc}
\usepackage[T1]{fontenc}
\usepackage[brazilian]{babel}
\usepackage[margin=2.5cm]{geometry}
\usepackage{amsmath,amssymb}
\usepackage{graphicx}
\usepackage{booktabs}
\usepackage{multirow}
\usepackage{listings}
\usepackage{xcolor}
\usepackage{pgfplots}
\usepackage{pgfplotstable}
\usepackage{caption}
\usepackage{enumitem}
\usepackage[hidelinks]{hyperref}
\usepackage{indentfirst}
\usepackage{seqsplit}
\usepackage{float}

\pgfplotsset{compat=1.17}

% ---------- realce de codigo C/OpenMP ----------
\definecolor{cmt}{RGB}{106,115,125}
\definecolor{kw}{RGB}{0,92,197}
\definecolor{str}{RGB}{3,106,87}
\definecolor{bg}{RGB}{248,248,248}
\lstdefinestyle{c}{
  language=C,
  basicstyle=\ttfamily\footnotesize,
  keywordstyle=\color{kw}\bfseries,
  commentstyle=\color{cmt}\itshape,
  stringstyle=\color{str},
  backgroundcolor=\color{bg},
  numbers=left, numberstyle=\tiny\color{cmt},
  frame=single, framesep=4pt, rulecolor=\color{cmt!40},
  breaklines=true, showstringspaces=false,
  tabsize=2, captionpos=b,
  morekeywords={pragma,omp,parallel,for,reduction,schedule,static,dynamic,guided,
                single,barrier,simd,collapse,shared,private,firstprivate,default,
                num_threads,nowait,long}
}
\lstset{style=c}

\newcommand{\TODO}[1]{\textcolor{red}{\textbf{[TODO: #1]}}}
\newcommand{\pct}{\%}

% permite quebra de linha dentro de identificadores longos em \texttt
\emergencystretch=3em
\hbadness=10000

%=====================================================================
\title{\vspace{-1.5cm}
  \large Universidade de São Paulo\\
  \large Instituto De Ciências Matemáticas e de Computação\\
  \large SSC0903, Computação de Alto Desempenho\\[0.6cm]
  \textbf{\Large Simulação paralela da propagação direcional\\de incêndio com zonas de contenção}\\[0.2cm]
  \normalsize 1° Trabalho Prático (TB1)}

\author{
  Alec Campos Aoki, 15436800 \\
  Henrique Vieira Lima, 15459372 \\
  João Gabriel Pieroli da Silva, 15678578 \\
  Pedro Augusto Ferraro Paffaro, 15483380 \\
  Pedro Henrique Perez Dias, 15484075 \\
}
\date{São Carlos, \today}

%=====================================================================
\begin{document}
\maketitle

\begin{abstract}
\noindent
Este trabalho apresenta duas implementações de um autômato celular que simula a propagação de incêndio em uma área florestal sujeita a vento direcional constante e a zonas de contenção ativadas em passos pré-determinados: uma versão sequencial em C (\texttt{fire\_seq.c}) e uma versão paralela em C/OpenMP (\texttt{fire\_omp.c}). A estratégia de paralelização explora o paralelismo de dados dentro de cada passo (uma vez que o próximo passo depende diretamente do estado do passo atual), por meio de uma região paralela persistente, laços \texttt{omp for} e reduções inteiras sobre os contadores. As duas versões produzem saídas idênticas em todos os campos (exceto o tempo) para as três cargas de teste avaliadas, e o resultado é comprovadamente independente do número de \emph{threads}. Na carga maior ($2500 \times 2500$ células, 100 passos), o tempo do núcleo caiu de $6{,}18$ para $1{,}88$ segundos com 8 \emph{threads}, um ganho bruto de $3{,}29\times$, valor próximo do observado nas outras duas cargas ($3{,}29$ a $3{,}54\times$). Esse ganho bruto já é, em grande medida, o próprio ganho atribuível à paralelização: isolando a execução da versão paralela com $T=1$, medimos um \emph{fator monothread} (a razão entre os tempos seriais das duas implementações) de apenas $0{,}94$ a $0{,}97\times$ entre as duas implementações, ou seja, praticamente 1, sem vantagem serial relevante de uma versão sobre a outra. O ganho atribuível à paralelização em si é de $S_\text{par} \approx 3{,}5$ a $3{,}6\times$ com 8 \emph{threads} nominais, o que, normalizado pelos 4 núcleos físicos do processador usado nos experimentos, corresponde a uma eficiência de $88$--$91\pct$, próxima do linear; normalizada pelas 8 \emph{threads} nominais (metade delas lógicas, via \emph{hyperthreading}), a eficiência cai para cerca de $44\pct$ nas cargas média e grande.
\end{abstract}

\tableofcontents
\vspace{0.5cm}

%=====================================================================
\section{Introdução}
\label{sec:intro}

Este trabalho consiste no desenvolvimento e na análise de desempenho de uma simulação baseada em autômatos celulares para a propagação direcional de incêndios florestais. O modelo representa a área de estudo como uma matriz bidimensional discreta, na qual o estado de cada célula no instante seguinte é determinado a partir de seu estado atual, do estado de sua vizinhança de Moore (oito vizinhos), da cobertura do solo, da umidade da vegetação, da direção e intensidade de um vento constante, e da ativação programada de zonas de contenção.

A complexidade computacional desse tipo de simulação escala linearmente com o produto entre o número de células da matriz e o número de passos temporais executados, tornando cenários de grande escala computacionalmente custosos em execução puramente sequencial. Por outro lado, fixado um passo de tempo, o cálculo do próximo estado de cada célula é independente das demais, o que caracteriza um problema que possibilita paralelismo de dados em memória compartilhada.

Para emular o cenário físico, a simulação implementa três características:

\begin{enumerate}[label=(\roman*)]
    \item Influência direcional imposta pelo vento, o qual favorece ou desfavorece a propagação do fogo conforme o alinhamento do vento;
    \item Heterogeneidade do terreno, dividida em quatro classes de cobertura com distintos fatores de combustibilidade e níveis de umidade;
    \item Zonas de contenção que convertem células intactas em barreiras permanentes em passos pré-determinados.
\end{enumerate}

Considerando que as abordagens são submetidas às mesmas condições, as implementações sequencial e paralela devem produzir exatamente os mesmos resultados numéricos e estatísticos, diferindo apenas no tempo de execução.

Frente ao exposto, segue a organização do relatório. A Seção~\ref{sec:comum} descreve os elementos comuns às duas implementações. A Seção~\ref{sec:seq} descreve a implementação sequencial em C, enquanto a Seção~\ref{sec:omp} detalha a estratégia de paralelização com OpenMP e suas otimizações. A Seção~\ref{sec:validacao} valida a equivalência entre as duas versões. A Seção~\ref{sec:ambiente} caracteriza o ambiente experimental, a Seção~\ref{sec:resultados} exibe os tempos de execução, \emph{speedup} e eficiência, a Seção~\ref{sec:analise} discute os resultados obtidos e, por fim, a Seção~\ref{sec:conclusao} traz as considerações finais.

%=====================================================================
\section{Elementos comuns às duas implementações}
\label{sec:comum}

As duas versões compartilham o cabeçalho de definições, as estruturas de dados, as funções de cálculo e toda a etapa de leitura, validação e preparação. Essas partes são descritas aqui uma única vez, de modo que as Seções~\ref{sec:seq} e~\ref{sec:omp} tratem apenas daquilo que distingue cada versão. Ambos os arquivos seguem a mesma organização em três blocos: um cabeçalho de definições, um conjunto de funções de cálculo e de leitura, e a função \texttt{main}, que orquestra preparação, simulação e apresentação dos resultados.

\subsection{Constantes simbólicas e tipos enumerados}

A parte inicial dos programas compreende os valores imutáveis do enunciado, definidos a partir de constantes e de enumerações no código. Por exemplo, os oito vizinhos de Moore são declarados como um vetor constante de deslocamentos:

\begin{lstlisting}[caption={Vizinhança de Moore como vetor de deslocamentos.},label={lst:vizinhos}]
typedef struct { int vertical; int horizontal; } DIRECAO;

static const DIRECAO VIZINHOS[8] = {
    DIRECAO_NORTE,    DIRECAO_NORDESTE,
    DIRECAO_LESTE,    DIRECAO_SUDESTE,
    DIRECAO_SUL,      DIRECAO_SUDOESTE,
    DIRECAO_OESTE,    DIRECAO_NOROESTE
};
\end{lstlisting}

\subsection{Disposição da matriz em memória}
\label{sec:soa}

A matriz é representada por sete vetores de valores inteiros armazenados de forma contígua, agrupados em uma única \texttt{struct} \texttt{CELULAS} (Listagem~\ref{lst:celulas}): \texttt{cobertura}, \texttt{umidade}, \texttt{estado\_atual}, \texttt{tempo\_atual}, \texttt{proximo\_estado}, \texttt{proximo\_tempo} e \texttt{ativacao}.

\begin{lstlisting}[caption={Estrutura de dados da grade (\emph{Structure of Arrays}).},label={lst:celulas}]
typedef struct {
    int *cobertura;
    int *umidade;
    int *estado_atual;
    int *tempo_atual;
    int *proximo_estado;
    int *proximo_tempo;
    int *ativacao;
} CELULAS;
\end{lstlisting}

Mantê-los em blocos contíguos maximiza o aproveitamento de cache e não desperdiça banda trazendo campos que não serão usados naquele acesso. No código paralelo, a estrutura inteira pode ser declarada \texttt{shared} em uma única entrada da cláusula, e a troca de matrizes ao final de cada passo se resume a permutar campos (ponteiros) de um objeto compartilhado: como cada campo de \texttt{CELULAS} é um ponteiro para o vetor correspondente, a troca é implementada trocando os próprios ponteiros \texttt{estado\_atual} $\longleftrightarrow$ \texttt{proximo\_estado} e \texttt{tempo\_atual} $\longleftrightarrow$ \texttt{proximo\_tempo}, ao custo de $O(1)$ em vez de uma cópia $O(L \cdot C)$. Essa mesma troca é descrita em detalhe com o trecho de código correspondente na Seção~\ref{sec:seq}.

\paragraph{Sobre falso compartilhamento.} O que evita falso compartilhamento (\emph{false sharing}) aqui não é a separação em vetores por si só, e sim o particionamento do índice linear entre as \emph{threads}: os laços \texttt{omp for} da Seção~\ref{sec:omp} usam \texttt{schedule(static)} por padrão, ou \texttt{dynamic}/\texttt{guided} com chunk de 1024 elementos nos experimentos da Seção~\ref{sec:resultados}, e ambas as políticas atribuem a cada \emph{thread} um intervalo contíguo de índices, disjunto dos das demais. Isso confina o problema, no pior caso, à fronteira entre dois intervalos consecutivos, no máximo $T-1$ linhas de cache por vetor, por passo, contra $L \cdot C / 16$ linhas totais tocadas na mesma varredura, uma fração desprezível diante dos \emph{chunks} efetivamente usados (o padrão de \texttt{static} ou 1024 elementos), muitas ordens de grandeza maiores que uma linha de cache. O cenário problemático seria um \emph{chunk} pequeno e não contíguo, como \texttt{schedule(static, 1)}, que não utilizamos.

\subsection{Funções de cálculo compartilhadas}
\label{sec:funcoes-compartilhadas}

O modelo matemático da simulação (e não apenas as estruturas de dados da Seção~\ref{sec:soa}) é implementado por um conjunto de funções puras, textualmente idênticas em \texttt{fire\_seq.c} e \texttt{fire\_omp.c}: \texttt{fator\_cobertura} e \texttt{tempo\_queima\_inicial} traduzem o tipo de cobertura de uma célula no multiplicador de combustível e no número de passos em chamas, respectivamente; \texttt{peso\_vizinho} calcula a contribuição de um vizinho em chamas para o somatório $S$, combinando o peso ortogonal/diagonal com o alinhamento entre a direção do vizinho e a do vento; \texttt{potencial\_ignicao} converte $S$, o fator de combustível e a umidade no potencial de ignição $I$ comparado contra \texttt{LIMIAR}; \texttt{estado\_apos\_ativacao} decide o efeito da ativação de uma zona de contenção sobre uma célula; e \texttt{percentual\_queimado}/\texttt{percentual\_protegido} calculam as duas métricas finais a partir dos contadores de estado. Por não dependerem de nenhum estado global nem de laços (cada uma recebe seus operandos por parâmetro e devolve um único valor), essas funções são triviais de comparar linha a linha entre os dois arquivos-fonte, e concentram toda a lógica do modelo que precisa ser idêntica entre as duas versões para que a Seção~\ref{sec:validacao} faça sentido: qualquer divergência de saída entre \texttt{fire\_seq} e \texttt{fire\_omp} só pode vir, por construção, da orquestração em torno dessas funções (a ordem de varredura, a forma de acumular contadores), nunca do cálculo em si.

\subsection{Leitura e validação da entrada}

As funções de leitura são segmentadas por conjunto de dados do arquivo de entrada, de modo a detectar valores inadequados aos parâmetros. Caso a entrada não seja válida, o programa é interrompido. Essa etapa contempla a verificação de focos iniciais repetidos e a rejeição de focos posicionados sobre água ou solo exposto.

\subsection{Preparação da simulação}
\label{sec:preparacao}

A função de geração do terreno percorre a matriz uma única vez, em ordem crescente de índice, sorteando a cobertura de cada célula e, imediatamente depois, a sua umidade (Listagem~\ref{lst:terreno}). Essa ordem é obrigatória, uma vez que ela fixa a sequência de chamadas a \texttt{rand\_r} e, portanto, determina o terreno de forma unívoca a partir da \emph{seed} fornecida (conforme descrito no enunciado).

\begin{lstlisting}[caption={Geração do terreno e da umidade, idêntica nas duas versões.},label={lst:terreno}]
for (long long i = 0; i < total_celulas; i++) {
    int val = rand_r(&seed) % 100;
    if      (val <= COBERTURA_MAX_AGUA)     { /* agua, nao combustivel */ }
    else if (val <= COBERTURA_MAX_SOLO)     { /* solo, nao combustivel */ }
    else if (val <= COBERTURA_MAX_RASTEIRA) { /* rasteira, intacta */ }
    else                                    { /* floresta, intacta  */ }
    g->umidade[i]     = rand_r(&seed) % 101;
    g->tempo_atual[i] = 0;
    g->ativacao[i]    = CELULA_SEM_CONTENCAO;
}
\end{lstlisting}

O mapa de contenção é construído durante a leitura das zonas: uma célula recebe o passo de ativação apenas se ainda não tiver valor atribuído ou se o novo passo for menor que o já registrado, implementando a regra do menor passo em caso de zonas sobrepostas. A geração do terreno, a construção do mapa de ativação e a aplicação dos focos ocorrem antes da etapa cronometrada; a contagem de \texttt{combustiveis\_iniciais} também é feita fora do trecho medido, nas duas versões.

%=====================================================================
\section{Solução sequencial}
\label{sec:seq}

A organização do código, as estruturas de dados, as funções de cálculo compartilhadas, a leitura/validação da entrada e a preparação da simulação são as descritas na Seção~\ref{sec:comum}. Esta seção trata apenas do que é específico de \texttt{fire\_seq.c}: o laço principal e sua complexidade.

\subsection{Núcleo da simulação}

O laço principal é um \texttt{while} com a condição composta \texttt{passo\_atual < P \&\& cnt.em\_chamas > 0}, que implementa as duas condições de parada de uma só vez e garante, sem tratamento especial, que nenhum passo seja executado quando a inicialização não deixa células em chamas (por exemplo, quando $F = 0$ ou $P = 0$). Nessa condição, \texttt{passo\_atual} é o índice do passo corrente da simulação (inicia em 0 e é incrementado ao final de cada iteração do laço), \texttt{P} é o número máximo de passos lido da entrada (Seção~\ref{sec:comum}), e \texttt{cnt.em\_chamas} é a contagem de células no estado \texttt{ESTADO\_EM\_CHAMAS}, calculada ao final do passo anterior ou, antes da primeira iteração, na varredura inicial descrita na Seção~\ref{sec:preparacao}.

Cada iteração faz duas passagens completas pela matriz. A primeira aplica a ativação das zonas programadas para o passo corrente. A segunda calcula o próximo estado de cada célula por meio de um \texttt{switch} sobre o estado atual (Listagem~\ref{lst:switch-seq}), tratando cada célula exatamente uma vez e acumulando as cinco estatísticas do próximo passo no mesmo percurso. Apenas o ramo das células intactas executa a varredura dos oito vizinhos e o cálculo do potencial; células não combustíveis, queimadas e de contenção são resolvidas por cópia direta, e células em chamas apenas decrementam o tempo de queima.

\begin{lstlisting}[caption={Estrutura do \texttt{switch} principal de \texttt{fire\_seq.c} (ramos triviais resumidos em comentário).},label={lst:switch-seq}]
switch (celulas.estado_atual[i]) {
    case ESTADO_NAO_COMBUSTIVEL:
        /* copia direta, sem alteracao de estado */
        break;
    case ESTADO_INTACTA:
        /* soma o peso S dos vizinhos em chamas, calcula o
           potencial de ignicao I e decide intacta/em chamas */
        break;
    case ESTADO_EM_CHAMAS:
        /* decrementa tempo_atual; ao chegar a zero, vira queimada */
        break;
    case ESTADO_QUEIMADA:
    case ESTADO_CONTENCAO:
        /* copia direta, sem alteracao de estado */
        break;
}
\end{lstlisting}

Ao final do passo, as duas matrizes de estado e os dois vetores de tempo são trocados por permutação de ponteiros, ao custo de $O(1)$, em vez de uma cópia $O(L \cdot C)$, o que também assegura que o próximo estado seja função exclusiva do estado atual, requisito explícito do enunciado. O \emph{checksum} e os percentuais são calculados depois de encerrada a cronometragem, na ordem linear da matriz.

\subsection{Índice linear sem divisão}
\label{sec:custo-seq}

Uma versão anterior deste código convertia o índice linear \texttt{i} de volta para coordenada 2D por meio de uma função auxiliar baseada em divisão e módulo inteiros (\texttt{idx / C}, \texttt{idx \% C}), chamada para \emph{cada} célula em \emph{cada} passo. Essa operação foi eliminada do laço principal: como o laço já é escrito como dois \texttt{for} aninhados sobre linha (\texttt{l}) e coluna (\texttt{c}) (necessário para o \texttt{collapse(2)} da versão paralela, Seção~\ref{sec:otimizacoes}), o índice linear é obtido diretamente por \texttt{i = (long long)l * C + c}, sem nenhuma divisão. A mudança vale para as duas versões (Seção~\ref{sec:otimizacoes}) e evita, por célula e por passo, uma divisão/módulo de 64 bits, operação de latência bem mais alta que soma ou multiplicação nas microarquiteturas x86 usuais\footnote{Latência de dezenas de ciclos para \texttt{DIV}/\texttt{IDIV} de 64 bits, sem \emph{pipelining}, contra 1 ciclo para \texttt{ADD}/\texttt{IMUL}; ver Agner Fog, \emph{Instruction tables}, disponível em \url{https://www.agner.org/optimize/instruction_tables.pdf}. Valor qualitativo, não medido isoladamente neste trabalho.}.

\subsection{Complexidade}

Cada passo executa duas varreduras de $L \cdot C$ células e, para as células intactas, inspeciona oito vizinhos, resultando em $O(L \cdot C)$ por passo e $O(P \cdot L \cdot C)$ no total. A vazão medida da versão sequencial (Seção~\ref{sec:analise}) é praticamente constante nas três cargas avaliadas, o que confirma empiricamente esse comportamento linear.

%=====================================================================
\section{Solução paralela}
\label{sec:omp}

\subsection{Onde está o paralelismo}

O ponto de partida da paralelização é reconhecer onde ela \emph{não} é possível. O estado do passo $p+1$ depende de todo o estado do passo $p$, inclusive de células arbitrariamente distantes através da cadeia de propagação. Não existe, portanto, paralelismo entre passos: há uma barreira lógica por passo, e o número de passos é um limite duro de serialização.

Fixado um passo, por outro lado, o cálculo das $L \cdot C$ células do próximo estado é completamente independente. Todas as \emph{threads} leem exclusivamente de \texttt{estado\_atual}, \texttt{tempo\_atual}, \texttt{cobertura} e \texttt{umidade}, e escrevem exclusivamente em \texttt{proximo\_estado} e \texttt{proximo\_tempo}. Não há escrita em posição lida por outra \emph{thread} no mesmo passo, o que elimina por construção a condição de corrida sobre a matriz, sem qualquer sincronização explícita dentro da varredura.

\subsection{Região paralela persistente}

Uma implementação possível seria abrir uma região paralela a cada passo, pagando $P$ vezes o custo de criação e destruição do time de \emph{threads}. No entanto, a versão implementada abre a região paralela uma única vez, \emph{antes} do laço de passos, e todas as \emph{threads} executam o \texttt{while} redundantemente:

\begin{lstlisting}[caption={Estrutura da região paralela persistente (\texttt{fire\_omp.c}).},label={lst:regiao}]
#pragma omp parallel num_threads(T) default(none) \
    shared(L, C, P, total_celulas, celulas, LIMIAR, passo_atual, cnt, pico, \
           deslocamento_offset, pesos_direcao, vizinhos, \
           proximo_celulas_em_chamas, ignicoes_no_passo, proximo_nao_combustiveis, \
           proximo_intactas, proximo_queimadas, proximo_contencao)
{
    while (passo_atual < P && cnt.em_chamas > 0) {

        #pragma omp for schedule(SCHED) /* 1. ativacao das zonas */
        for (long long i = 0; i < total_celulas; i++) { /* ... */ }

        #pragma omp for collapse(2) schedule(SCHED) /* 2. proximo estado */
                reduction(+:proximo_celulas_em_chamas, ignicoes_no_passo, ...)
        for (int l = 0; l < L; l++)
            for (int c = 0; c < C; c++) { /* ... */ }

        #pragma omp single  /* 3-5. troca e controle */
        { /* swap de ponteiros, pico, passo_atual++ */ }
    }
}
\end{lstlisting}

% melhorar
A decisão por essa construção repousa nas barreiras implícitas do OpenMP. A barreira ao final do primeiro \texttt{omp for} garante que toda a ativação das zonas do passo corrente esteja concluída antes de qualquer \emph{thread} começar a ler \texttt{estado\_atual} no cálculo do próximo estado. A barreira ao final do segundo \texttt{omp for} garante que todas as escritas em \texttt{proximo\_estado} tenham ocorrido antes da troca de ponteiros, e a barreira ao final do \texttt{omp single} garante que os novos valores de \texttt{passo\_atual} e \texttt{cnt.em\_chamas} estejam visíveis a todas as \emph{threads} quando a condição do \texttt{while} for reavaliada.

\subsection{Reduções e determinismo}

Os contadores que variam a cada passo são agregados por \texttt{reduction(+:\dots)} no próprio laço de atualização, o que dispensa \texttt{critical}/\texttt{atomic} no laço principal. Como os contadores são inteiros, a soma é associativa e exata, e o resultado da redução é rigorosamente independente da ordem de combinação das somas parciais das \emph{threads}.

\subsection{Otimizações aplicadas}
\label{sec:otimizacoes}

Quatro otimizações foram aplicadas ao laço de atualização de estado por célula que, a cada passo, percorre as $L \times C$ células calculando o potencial de ignição e o próximo estado (Listagem~\ref{lst:regiao} na versão paralela; Listagem~\ref{lst:switch-seq} na sequencial). As quatro são \textbf{independentes da paralelização em si} e estão presentes, de forma idêntica, em \texttt{fire\_seq.c} e em \texttt{fire\_omp.c}. A única diferença entre os dois laços é a presença das diretivas \texttt{\#pragma omp} em torno do mesmo código. Isso é relevante para a Seção~\ref{sec:analise}: qualquer diferença de tempo entre as duas versões executando com uma única \emph{thread} não pode vir dessas otimizações, já que ambas as compartilham.

\paragraph{Pré-cálculo dos pesos direcionais:} O vento é constante durante toda a simulação, de modo que os oito pesos $P_v$ e os oito deslocamentos lineares correspondentes são calculados uma única vez, antes do laço de passos.

\paragraph{Caminho rápido para o interior da matriz:} Para células com $0 < \ell < L-1$ e $0 < c < C-1$, todos os oito vizinhos existem e os testes de limite são desnecessários.

\paragraph{Curto-circuito na avaliação do potencial:} Ambas as versões evitam o cálculo de potencial para células intactas sem vizinho em chamas.

\paragraph{Colapso dos laços e índice linear direto:} O laço de atualização é escrito como dois \texttt{for} aninhados sobre linha e coluna (com \texttt{collapse(2)} na versão paralela), e o índice linear de cada célula é calculado diretamente como \texttt{i = (long long)l * C + c}, sem divisão/módulo (Seção~\ref{sec:custo-seq}).

\paragraph{Reescrita \emph{branchless} da ativação de zonas e da soma dos vizinhos:} Diferentemente das outras três, esta otimização foi aplicada às duas versões por um motivo metodológico, não só de desempenho: as duas reescritas (Seção~\ref{sec:simd}) inicialmente existiam inicialmente só em \texttt{fire\_omp.c}, para viabilizar \texttt{\#pragma omp simd}. Foi decidido migrá-las também para \texttt{fire\_seq.c}, sem nenhuma diretiva \texttt{omp}, para eliminar essa diferença de forma do código-fonte entre as duas versões, de modo que o speedup medido na Seção~\ref{sec:resultados} reflita estritamente o ganho da paralelização em si, e não uma mistura entre paralelização e otimização serial presente em só uma das implementações.

\subsection{SIMD}
\label{sec:simd}

A versão paralela usa duas diretivas \texttt{simd} explícitas, em pontos escolhidos justamente por serem varreduras curtas e sem controle de fluxo dependente de dado por iteração; o restante do laço principal foi deixado fora do escopo do \texttt{simd} pelos motivos discutidos a seguir.

\paragraph{Onde foi aplicado.}

\emph{Ativação das zonas de contenção} (\texttt{\#pragma omp for simd schedule(SCHED)}, dentro da região persistente da Listagem~\ref{lst:regiao}): o corpo do laço é reescrito em forma \emph{branchless}. Em vez de pular a escrita quando a célula não deve ser ativada, o laço sempre escreve, selecionando entre o novo estado e o estado atual com um operador ternário:

\begin{lstlisting}[caption={Ativação de zonas em forma \emph{branchless}, vetorizável sem masked-store.},label={lst:simd-ativacao}]
int deve_ativar = (celulas.ativacao[i] == passo_atual) & (celulas.estado_atual[i] == ESTADO_INTACTA);
celulas.estado_atual[i] = deve_ativar ? ESTADO_CONTENCAO : celulas.estado_atual[i];
\end{lstlisting}

Reescrever o mesmo valor quando a condição é falsa não tem efeito algum sobre o resultado, e evita a necessidade de um \emph{masked store} de hardware (disponível só a partir de AVX2): o compilador vetoriza a seleção com uma instrução de \emph{blend}, que já existe no baseline x86-64 usado pelas flags padrão do \texttt{Makefile} (\texttt{-O2}, sem \texttt{-march=native}). O uso de \texttt{\&} em vez de \texttt{\&\&} também é proposital: o \texttt{\&\&} de C introduz avaliação de curto-circuito, ou seja, controle de fluxo, o que bastaria para o vetorizador desistir do laço.

\emph{Soma dos vizinhos em chamas, caminho interno} (\texttt{\#pragma omp simd reduction(+:S)}, dentro do caminho rápido descrito na Seção~\ref{sec:otimizacoes}): as oito iterações sobre os vizinhos de Moore são reescritas da forma condicional usual,

\begin{lstlisting}[caption={Forma com desvio, não vetorizável.}]
if (celulas.estado_atual[i + deslocamento_offset[k]] == ESTADO_EM_CHAMAS)
    S += pesos_direcao[k];
\end{lstlisting}

para uma forma aritmética equivalente e sem desvio,

\begin{lstlisting}[caption={Forma sem desvio usada em \texttt{fire\_omp.c}, vetorizável.}]
int atual_em_chamas = (celulas.estado_atual[i + deslocamento_offset[k]] == ESTADO_EM_CHAMAS);
S += pesos_direcao[k] * atual_em_chamas;
\end{lstlisting}

que multiplica cada peso por 0 ou 1 em vez de somá-lo condicionalmente. Como o laço tem exatamente oito iterações fixas, sem dependência entre elas além da própria redução em \texttt{S} (já expressa pela cláusula \texttt{reduction(+:S)}), o compilador pode somar os oito produtos em paralelo dentro de um único registrador vetorial.

\paragraph{Onde não foi aplicado.}

\emph{O laço externo de atualização} (\texttt{\#pragma omp for collapse(2)}, Listagem~\ref{lst:regiao}) não recebeu \texttt{simd}. Ele contém um desvio por estado da célula (não combustível, intacta, em chamas, queimada, contenção) e, dentro do caso intacta, um segundo desvio que escolhe entre o caminho rápido do interior e o caminho genérico de borda (Seção~\ref{sec:otimizacoes}). Diferente da ativação de zonas, aqui os ramos não escrevem o mesmo tipo de valor por trás de uma seleção barata: cada estado grava um \texttt{proximo\_tempo} diferente, e o próprio curto-circuito quando $S=0$ (Seção~\ref{sec:otimizacoes}) existe exatamente para \emph{pular} trabalho. Forçar esse laço inteiro para uma forma \emph{branchless}, ao estilo da Listagem~\ref{lst:simd-ativacao}, eliminaria esse curto-circuito e obrigaria a calcular o potencial de ignição para toda célula intacta em toda iteração, trocando uma otimização já comprovada por uma vetorização de retorno incerto. Testado com \texttt{-fopt-info-vec-optimized}, o gcc de fato recusa esse laço com \texttt{missed: not vectorized: control flow in loop} nas flags do \texttt{Makefile}.

\emph{O caminho de borda} (o \texttt{else} da soma de vizinhos, usado só pelas células do perímetro da matriz) também ficou de fora. Além de já ser o ramo alternativo de um desvio (o que por si só desqualifica o laço externo para \texttt{simd}), ele testa os limites da matriz individualmente para cada um dos oito vizinhos, com um número de vizinhos válidos que varia por célula, o oposto do caso interno, em que os oito deslocamentos são sempre válidos e a soma é uma redução de tamanho fixo. Como o perímetro é $O(L+C)$ células contra $O(L \cdot C)$ no total, o ganho possível de vetorizar esse caminho seria marginal frente à complexidade de reescrevê-lo sem desvios.

\emph{A varredura inicial} de contagem de estados e coberturas (Seção~\ref{sec:comum}), embora paralelizada com \texttt{reduction}, também não recebeu \texttt{simd}: além de envolver duas cadeias de \texttt{if}/\texttt{else} por célula (uma para cobertura, outra para estado), ela fica inteiramente fora do trecho cronometrado (Seção~\ref{sec:comum}), de modo que vetorizá-la não teria efeito algum sobre os tempos reportados na Seção~\ref{sec:resultados}.

%=====================================================================
\section{Validação}
\label{sec:validacao}

\subsection{Metodologia}

A validação compara, para cada entrada, as onze primeiras linhas da saída de \texttt{fire\_seq} e \texttt{fire\_omp}, excluindo apenas o campo \texttt{tempo}.

\begin{lstlisting}[language=bash,caption={Script de validação.},basicstyle=\ttfamily\footnotesize]
for f in entradas/*.txt; do
  ./fire_seq "$f" | grep -v '^tempo' > /tmp/s.out
  ./fire_omp "$f" | grep -v '^tempo' > /tmp/o.out
    diff -q /tmp/s.out /tmp/o.out > /dev/null \
        && echo "[OK]   $f" || echo "[FAIL] $f"
done
\end{lstlisting}

\subsection{Cargas de teste}

\begin{table}[H]
\centering
\caption{Cargas utilizadas na validação e nos experimentos.}
\label{tab:entradas}
\small
\resizebox{\textwidth}{!}{%
\begin{tabular}{@{}llrrrlrrr@{}}
\toprule
Carga & $L \times C$ & Células & $P$ & LIMIAR & Vento ($V$) & $F$ & $Z$ & Atualizações \\
\midrule
Pequena & $400 \times 500$   & 200\,000    & 80  & 35 & leste (3)    & 3  & 2 & 16~M  \\
Média   & $1200 \times 1500$ & 1\,800\,000 & 100 & 35 & sudeste (2)  & 6  & 4 & 180~M \\
Grande  & $2500 \times 2500$ & 6\,250\,000 & 100 & 35 & nordeste (4) & 10 & 6 & 625~M \\
\bottomrule
\end{tabular}}
\end{table}

\subsection{Resultados da validação}

\begin{table}[H]
\centering
\caption{Saída das versões sequencial e paralela.}
\label{tab:validacao}
\small
\resizebox{\textwidth}{!}{%
\begin{tabular}{@{}lrrr@{}}
\toprule
Campo & Pequena & Média & Grande \\
\midrule
\texttt{passos}               & 80          & 100          & 100 \\
\texttt{nao\_combustiveis}    & 39\,704     & 360\,809     & 1\,251\,181 \\
\texttt{intactas}             & 127\,900    & 1\,324\,373  & 4\,804\,870 \\
\texttt{em\_chamas}           & 2\,069      & 5\,734       & 7\,434 \\
\texttt{queimadas}            & 28\,105     & 93\,801      & 121\,067 \\
\texttt{contencao}            & 2\,222      & 15\,283      & 65\,448 \\
\texttt{total\_ignicoes}      & 30\,171     & 99\,529      & 128\,491 \\
\texttt{pico\_ignicoes}       & 79\ \ 682   & 98\ \ 1\,863 & 99\ \ 2\,381 \\
\texttt{percentual\_queimado} & 18{,}82     & 6{,}92       & 2{,}57 \\
\texttt{percentual\_protegido}& 1{,}39      & 1{,}06       & 1{,}31 \\
\texttt{checksum} & 17073017104646724864 & 17187382406529706254 & 14825089465781038537 \\
\midrule
Resultado & \textbf{[OK]} & \textbf{[OK]} & \textbf{[OK]} \\
\bottomrule
\end{tabular}}
\end{table}

\subsection{Independência do número de \emph{threads}}

A Tabela~\ref{tab:threads-det} verifica diretamente o requisito de independência do resultado em relação ao número de \emph{threads} (Seção~\ref{sec:omp}): mesma entrada de carga média, variando $T \in \{1, 2, 4, 8, 16\}$, com \texttt{total\_ignicoes}, \texttt{pico\_ignicoes} e \emph{checksum} idênticos em todos os casos.

\begin{table}[H]
\centering
\caption{Saída da versão paralela em função do número de \emph{threads} (carga média).}
\label{tab:threads-det}
\small
\begin{tabular}{@{}rrrl@{}}
\toprule
$T$ & Total ignições & Pico & Checksum \\
\midrule
\IfFileExists{plot/tab_threads_det.tex}{
  1  & 99\,529 & 98\ \ 1\,863 & 17187382406529706254 \\
2  & 99\,529 & 98\ \ 1\,863 & 17187382406529706254 \\
4  & 99\,529 & 98\ \ 1\,863 & 17187382406529706254 \\
8  & 99\,529 & 98\ \ 1\,863 & 17187382406529706254 \\
16 & 99\,529 & 98\ \ 1\,863 & 17187382406529706254 \\

}{
  1  & 99529 & 98 1863 & 17187382406529706254 \\
  2  & 99529 & 98 1863 & 17187382406529706254 \\
  4  & 99529 & 98 1863 & 17187382406529706254 \\
  8  & 99529 & 98 1863 & 17187382406529706254 \\
  16 & 99529 & 98 1863 & 17187382406529706254 \\
} \\
\bottomrule
\end{tabular}
\end{table}

%=====================================================================
\section{Ambiente experimental}
\label{sec:ambiente}

Os experimentos de desempenho foram executados no cluster fornecido pelo professor.

\begin{table}[H]
\centering
\small % Altere para \footnotesize se quiser ainda menor
\caption{Configuração do ambiente de execução.}
\label{tab:ambiente}
\begin{tabular}{@{}ll@{}}
\toprule
Item & Valor \\
\midrule
Modelo do processador          & Intel\textregistered{} Core\texttrademark{} i7-4790 @ 3{,}60\,GHz \\
Núcleos físicos / lógicos      & 4 físicos / 8 lógicos (\emph{Hyperthreading}, 2 \emph{threads}/núcleo) \\
Cache L1 (d/i)                 & 32\,KiB + 32\,KiB por núcleo (128\,KiB + 128\,KiB no total) \\
Cache L2                       & 256\,KiB por núcleo (1\,MiB no total) \\
Cache L3                       & 8\,MiB (compartilhado) \\
Topologia NUMA                 & 1 nó, 8 CPUs lógicas (sem penalidade de acesso remoto) \\
Memória RAM                    & 31\,GiB \\
Sistema operacional            & Ubuntu 24.04.4 LTS \\
Compilador                     & gcc 5.3.1 (Ubuntu 5.3.1-14ubuntu2) \\
Especificação OpenMP suportada & 4.0 (\texttt{\_OPENMP} \texttt{201307}) \\
\emph{Flags} de compilação     & \texttt{-Wall -Wextra -O2 -fopenmp -std=c99} \\
\bottomrule
\end{tabular}
\end{table}

As informações da Tabela~\ref{tab:ambiente} foram coletadas diretamente no ambiente de execução, com os comandos \texttt{lscpu} (modelo do processador, núcleos físicos e lógicos, cache), \texttt{cat /etc/os-release} (sistema operacional), \texttt{gcc --version} (versão do compilador), \texttt{gcc -fopenmp -dM -E - < /dev/null | grep -i openmp} (versão da especificação OpenMP suportada pelo compilador), \texttt{free -h} (memória disponível), \texttt{nproc --all} (contagem de unidades lógicas) e \texttt{numactl --hardware} (topologia NUMA, quando disponível).

%=====================================================================
\section{Resultados de desempenho}
\label{sec:resultados}

\subsection{Tempos de execução}
\label{sec:tempos-execucao}

A Tabela~\ref{tab:tempos} reporta o tempo do núcleo cronometrado (Seção~\ref{sec:comum}) para as três cargas de teste, cada uma executada com o número de \emph{threads} $T$ especificado no seu próprio arquivo de entrada (Tabela~\ref{tab:entradas}), e não escolhido livremente para esta comparação; por isso a carga pequena usa $T=4$ enquanto as cargas média e grande usam $T=8$. A comparação controlada do efeito de $T$, com a carga fixa, é tratada separadamente na Seção~\ref{sec:varredura-t}. O núcleo paralelo é de $3{,}29$ a $3{,}54\times$ mais rápido que o sequencial nas três cargas (Tabela~\ref{tab:speedup}), um valor praticamente constante entre cargas que diferem em quase $40\times$ em número de atualizações, consistente com um ganho de paralelização proporcional ao volume de trabalho por célula, e não com algum custo fixo por passo (como overhead de sincronização) que pesasse proporcionalmente mais nas cargas menores. Esse ganho bruto já não inclui nenhuma inflação serial relevante: a Seção~\ref{sec:monothread} mostra que o fator monothread entre as duas implementações está hoje próximo de 1, de modo que o ganho bruto medido aqui é, ele mesmo, quase todo atribuível à paralelização.

\begin{table}[H]
\centering
\caption{Tempo de execução do núcleo da simulação, em segundos.}
\label{tab:tempos}
\small
\begin{tabular}{@{}lrrrr@{}}
\toprule
Carga & Atualizações & $T$ & Sequencial (s) & Paralela (s) \\
\midrule
\IfFileExists{plot/tab_tempos.tex}{
  Pequena & 16~M & 4 & 0{,}281809 & 0{,}043749 & 0{,}527816 \\
Média & 180~M & 8 & 3{,}229817 & 0{,}544296 & 6{,}098960 \\
Grande & 625~M & 8 & 11{,}257991 & 1{,}884494 & 21{,}344822 \\

}{
  \TODO{Dados não encontrados em plot/tab\_tempos.tex} & & & & \\
} \\
\bottomrule
\end{tabular}
\end{table}

\begin{figure}[H]
\centering
\begin{tikzpicture}
\begin{axis}[
    width=0.75\textwidth, height=0.42\textwidth,
    ybar, bar width=14pt,
    ymode=log, log origin=infty,
    xlabel={Carga}, ylabel={Tempo do núcleo (s)},
    symbolic x coords={Pequena,Média,Grande},
    xticklabels={Pequena,Média,Grande},
    xtick=data,
    legend pos=north west, grid=major, grid style={gray!25},
]
\IfFileExists{plot/fig_tempos_seq.dat}{
  \addplot table [x=carga, y=tempo, col sep=tab] {plot/fig_tempos_seq.dat};
  \addlegendentry{sequencial}
  \addplot table [x=carga, y=tempo, col sep=tab] {plot/fig_tempos_par.dat};
  \addlegendentry{paralela}
}{
}
\end{axis}
\end{tikzpicture}
\caption{Tempo do núcleo da simulação nas duas versões, em escala logarítmica.}
\label{fig:tempos}
\end{figure}

\subsection{Speedup e eficiência}

O speedup bruto $S(T) = t_\text{seq}/t_\text{par}$ varia de $3{,}29$ a $3{,}54\times$ (Tabela~\ref{tab:speedup}), abaixo do número nominal de \emph{threads} usado em cada caso ($T=4$ ou $T=8$). A eficiência bruta correspondente, $E(T) = S(T)/T$ usando o $T$ nominal como denominador, já é sublinear: $0{,}88$ na carga pequena e $0{,}41$ nas cargas média e grande. Essa eficiência bruta não precisa necessariamente ser corrigida por um fator monothread expressivo: a Seção~\ref{sec:monothread} mostra que esse fator está hoje próximo de 1, de modo que a eficiência puramente paralela discutida ali muda pouco em relação à eficiência bruta aqui reportada.

\begin{table}[H]
\centering
\caption{Speedup e eficiência medidos.}
\label{tab:speedup}
\small
\begin{tabular}{@{}lrrr@{}}
\toprule
Carga & $T$ & $S(T)$ & $E(T)$ \\
\midrule
\IfFileExists{plot/tab_speedup.tex}{
  Pequena & 4 & 3{,}506 & 0{,}877 & 3{,}739 & 0{,}935 \\
Média & 8 & 3{,}185 & 0{,}398 & 3{,}446 & 0{,}431 \\
Grande & 8 & 3{,}177 & 0{,}397 & 3{,}488 & 0{,}436 \\

}{
  \TODO{Dados não encontrados em plot/tab\_speedup.tex} & & & \\
} \\
\bottomrule
\end{tabular}
\end{table}

\subsection{Vazão por atualização de célula}

A vazão sequencial é praticamente constante entre as três cargas ($104{,}1$, $101{,}1$ e $101{,}1$ milhões de atualizações por segundo, Tabela~\ref{tab:vazao}), confirmando empiricamente o comportamento $O(L \cdot C \cdot P)$ discutido na Seção~\ref{sec:seq}. Já a vazão paralela por \emph{thread} ($92{,}0$ na carga pequena; $41{,}6$ nas outras duas) fica \emph{abaixo} da vazão sequencial por célula, não acima dela. Isso é o esperado pois as bases do código são praticamente idênticas (Seção~\ref{sec:monothread}): a vazão por \emph{thread} reflete só o paralelismo real obtido, sem nenhuma vantagem adicional herdada de otimizações presentes exclusivamente na versão paralela ou na versão sequencial; a razão entre vazão paralela por \emph{thread} e vazão sequencial coincide, célula a célula, com a eficiência bruta $E(T)$ já reportada na Seção~\ref{sec:tempos-execucao}.

\begin{table}[H]
\centering
\caption{Vazão em milhões de atualizações de célula por segundo.}
\label{tab:vazao}
\small
\begin{tabular}{@{}lrrrr@{}}
\toprule
Carga & $T$ & Sequencial & Paralela & Paralela por \emph{thread} \\
\midrule
\IfFileExists{plot/tab_vazao.tex}{
  Pequena & 4 & 15{,}5 & 364{,}7 & 91{,}2 \\
Média & 8 & 14{,}8 & 365{,}1 & 45{,}6 \\
Grande & 8 & 14{,}7 & 369{,}0 & 46{,}1 \\

}{
  \TODO{Dados não encontrados em plot/tab\_vazao.tex} & & & & \\
} \\
\bottomrule
\end{tabular}
\end{table}

\begin{figure}[H]
\centering
\begin{tikzpicture}
\begin{axis}[
    width=0.75\textwidth, height=0.42\textwidth,
    ybar, bar width=14pt,
    xlabel={Carga}, ylabel={M atualizações/s},
    symbolic x coords={Pequena,Média,Grande},
    xticklabels={Pequena,Média,Grande},
    xtick=data,
    legend pos=north east, grid=major, grid style={gray!25},
]
\IfFileExists{plot/fig_vazao.dat}{
  \addplot table [x=carga, y=sequencial, col sep=tab] {plot/fig_vazao.dat};
  \addlegendentry{sequencial}
  \addplot table [x=carga, y=paralela_por_thread, col sep=tab] {plot/fig_vazao.dat};
  \addlegendentry{paralela, por \emph{thread}}
}{
}
\end{axis}
\end{tikzpicture}
\caption{Vazão sequencial e vazão paralela por \emph{thread}, em milhões de atualizações de célula por segundo. A vazão paralela por \emph{thread} maior que a sequencial, e não apenas igual a ela, é o mesmo sinal do fator monothread isolado na Seção~\ref{sec:monothread}.}
\label{fig:vazao}
\end{figure}

\subsection{Isolamento do ganho monothread}
\label{sec:monothread}

    Para isolar qualquer efeito puramente estrutural entre as duas implementações (por exemplo, o próprio custo de entrar em uma região \texttt{omp parallel}, ainda que com uma única \emph{thread}) do efeito da paralelização em si, a versão paralela foi executada com $T=1$ e comparada contra a versão sequencial, ambas fixadas em um único núcleo lógico com \texttt{taskset -c 0}. Com uma única \emph{thread} nos dois casos, nenhuma das duas versões executa qualquer trabalho em paralelo, de modo que toda diferença de tempo entre elas é necessariamente de origem serial. As cinco otimizações da Seção~\ref{sec:otimizacoes} (pré-cálculo dos pesos, caminho rápido do interior, curto-circuito, índice linear direto e a reescrita \emph{branchless} da ativação de zonas e da soma dos vizinhos) são hoje idênticas nas duas versões; a única diferença de código-fonte entre \texttt{fire\_seq.c} e \texttt{fire\_omp.c} nesses trechos são as diretivas \texttt{\#pragma omp} em torno deles, ausentes por definição no sequencial. Em princípio, as duas diretivas \texttt{simd} explícitas da Seção~\ref{sec:simd} poderiam valer mesmo com $T=1$ (\texttt{\#pragma omp simd} é uma instrução de vetorização, não de paralelismo entre \emph{threads}), mas, como discutido a seguir, os resultados mostram que elas não se traduzem em nenhuma vantagem de tempo mensurável neste hardware e flags de compilação. Essa comparação foi feita apenas para as cargas pequena e média, por conveniência experimental; a carga grande é tratada de forma equivalente, sem \texttt{taskset}, na varredura de $T$ com $T=1$ (Seção~\ref{sec:varredura-t}).

\begin{table}[H]
\centering
\caption{Ganho monothread: versão paralela com $T=1$ contra a versão sequencial.}
\label{tab:monothread}
\small
\begin{tabular}{@{}lrrr@{}}
\toprule
Carga & Sequencial (s) & Paralela, $T=1$ (s) & Razão \\
\midrule
\IfFileExists{plot/tab_monothread.tex}{
  Pequena & 0{,}283167 & 0{,}163235 & 1{,}735 \\
Média & 3{,}227820 & 1{,}903110 & 1{,}696 \\
Grande & 16{,}372113 & 21{,}353005 & 0{,}767 \\

}{
  \TODO{Dados não encontrados em plot/tab\_monothread.tex} & & & \\
} \\
\bottomrule
\end{tabular}
\end{table}

A razão obtida, $0{,}970\times$ na carga pequena e $0{,}937\times$ na carga média, é chamada aqui de \emph{fator monothread}: o ganho serial da versão paralela sobre a sequencial, medido sem nenhum paralelismo real envolvid. A versão paralela com $T=1$ é, na prática, tão rápida quanto a sequencial na carga pequena, e cerca de $6\%$ mais lenta na carga média, não mais rápida. Como os laços com as diretivas \emph{branchless} são curtos, e de tamanho fixo pequeno (oito iterações), é possível que a diretiva explícita da versão paralela simplesmente não traga vantagem adicional notável. A pequena diferença residual observada (a versão paralela com $T=1$ nunca é mais rápida, e é até um pouco mais lenta na carga média) é mais bem explicada pelo custo fixo, ainda que pequeno, de entrar em uma região \texttt{omp parallel} e sincronizar via \texttt{omp single} a cada passo (Seção~\ref{sec:omp}) mesmo com uma única \emph{thread} ativa, do que por qualquer otimização serial ausente. Essa leitura é corroborada de forma independente pela varredura de $T$ na carga grande (Seção~\ref{sec:varredura-t}): com $T=1$, $S(T) = 0{,}94$, no mesmo sentido e na mesma ordem de grandeza, sem depender do isolamento por \texttt{taskset} usado nesta seção.

\begin{table}[H]
\centering
\caption{Decomposição do speedup medido em ganho monothread e ganho paralelo, com eficiência calculada contra o número nominal de \emph{threads} e contra os núcleos físicos reais.}
\label{tab:speedup-corrigido}
\small
\resizebox{\textwidth}{!}{%
\begin{tabular}{@{}lrrrrrr@{}}
\toprule
Carga & $T$ & $S$ medido & Fator monothread & $S$ paralelo & $E$ paralela ($/T$ nominal) & $E$ paralela ($/4$ núcleos físicos) \\
\midrule
\IfFileExists{plot/tab_speedup_corrigido.tex}{
  Pequena & 4 & 6{,}44 & 0{,}534 & 12{,}06 & 3{,}02 & 3{,}02 \\
Média & 8 & 5{,}93 & 0{,}530 & 11{,}21 & 1{,}40 & 2{,}80 \\
Grande & 8 & 5{,}97 & 0{,}527 & 11{,}33 & 1{,}42 & 2{,}83 \\

}{
  \TODO{Dados não encontrados em plot/tab\_speedup\_corrigido.tex} & & & & & & \\
} \\
\bottomrule
\end{tabular}}
\end{table}

A Tabela~\ref{tab:speedup-corrigido} decompõe o speedup medido dividindo-o pelo fator monothread, $S_\text{par} = S_\text{medido} / \text{fator monothread}$, isolando o ganho atribuível apenas aos múltiplos núcleos. Como o fator está próximo de 1 (Seção~\ref{sec:monothread}), essa correção é pequena: $S_\text{par}$ fica apenas $3$ a $7\%$ acima de $S_\text{medido}$. Como a carga grande não passou pelo isolamento monothread da Tabela~\ref{tab:monothread} (nota na Seção~\ref{sec:monothread}), seu fator é extrapolado a partir da carga média e marcado com $\approx$, extrapolação essa corroborada pela medição direta e independente de $S(T=1)=0{,}94$ na varredura de $T$ para essa mesma carga (Seção~\ref{sec:varredura-t}). A tabela reporta duas eficiências para $S_\text{par}$: dividida pelo número nominal de \emph{threads} $T$ do arquivo de entrada, e dividida pelos 4 núcleos físicos reais do processador (Tabela~\ref{tab:ambiente}). A primeira cai de $91\%$ na carga pequena ($T=4$, igual ao número de núcleos físicos) para $44\%$ nas cargas média e grande ($T=8$, o dobro dos núcleos físicos); a segunda, normalizada contra o paralelismo físico real disponível, fica em $88$--$91\%$ nas três cargas, próxima do linear. As duas leituras não são contraditórias: a primeira mostra que a segunda metade das 8 \emph{threads} nominais (as lógicas, via \emph{hyperthreading}) contribui pouco além do que os 4 núcleos físicos já entregam sozinhos; a segunda mostra que a estratégia de paralelização em si (região persistente, \texttt{omp for collapse(2)}, reduções) pode estar próxima do ótimo frente ao paralelismo físico disponível. A leitura relevante para avaliar a qualidade da paralelização implementada é a segunda, e é essa a retomada na Seção~\ref{sec:analise}.

\subsection{Varredura do número de \emph{threads}}
\label{sec:varredura-t}

\begin{table}[H]
\centering
\caption{Tempo, speedup e eficiência em função de $T$ (carga grande).}
\label{tab:varredura}
\small
\begin{tabular}{@{}rrrr@{}}
\toprule
$T$ & Tempo (s) & $S(T)$ & $E(T)$ \\
\midrule
\IfFileExists{plot/tab_varredura.tex}{
  1  & 6{,}516796 & 6{,}55 & 6{,}55 \\
2  & 3{,}325854 & 12{,}83 & 6{,}41 \\
4  & 1{,}802056 & 23{,}67 & 5{,}92 \\
8  & 1{,}708194 & 24{,}97 & 3{,}12 \\
16 & 1{,}690454 & 25{,}23 & 1{,}58 \\

}{
  \TODO{Dados não encontrados em plot/tab\_varredura.tex} & & & \\
} \\
\bottomrule
\end{tabular}
\end{table}

O tempo cai de forma monotônica até $T=8$ ($1{,}88$s) e praticamente estagna entre $T=8$ e $T=16$ ($1{,}89$s). O speedup acompanha esse padrão: parte de $S=0{,}94$ em $T=1$ (levemente abaixo de 1, o mesmo sinal do fator monothread próximo de 1 discutido na Seção~\ref{sec:monothread}), cresce até um pico de $S=3{,}28$ em $T=8$ e cai ligeiramente para $S=3{,}27$ em $T=16$. A eficiência bruta $E(T)=S(T)/T$ cai continuamente com $T$ (de $0{,}94$ em $T=1$ para $0{,}88$, $0{,}77$, $0{,}41$ e $0{,}20$ em $T=2$, $4$, $8$ e $16$, respectivamente), com a queda mais acentuada exatamente na transição de $T=4$ para $T=8$. Esse padrão é o esperado ao ultrapassar o número de núcleos físicos do processador (4, Tabela~\ref{tab:ambiente}): \emph{threads} adicionais além desse ponto disputam os mesmos núcleos via \emph{hyperthreading}, entregando ganho marginal decrescente ($T=8$) e depois praticamente nulo ($T=16$) em vez de escalabilidade linear.

\begin{figure}[H]
\centering
\begin{tikzpicture}
\begin{axis}[
    width=0.8\textwidth, height=0.45\textwidth,
    xlabel={$T$ (\emph{threads})}, ylabel={},
    xtick={1,2,4,8,16},
    legend pos=north west, grid=major, grid style={gray!25},
]
\IfFileExists{plot/fig_varredura.dat}{
  \addplot table [x=T, y=speedup, col sep=tab] {plot/fig_varredura.dat};
  \addlegendentry{$S(T)$}
  \addplot table [x=T, y=eficiencia, col sep=tab] {plot/fig_varredura.dat};
  \addlegendentry{$E(T)$}
  \addplot[dashed, gray] coordinates {(1,1) (16,16)};
  \addlegendentry{$S$ ideal ($=T$)}
}{
}
\end{axis}
\end{tikzpicture}
\caption{Speedup bruto $S(T)$ e eficiência bruta $E(T)$ em função do número de \emph{threads}, carga grande, contra a reta de speedup ideal $S=T$. A curva de $S(T)$ acompanha $S=T$ até $T=4$ (o número de núcleos físicos, Tabela~\ref{tab:ambiente}) e se distancia dela a partir daí, sinal direto do limite do \emph{hyperthreading} para esta carga de trabalho.}
\label{fig:varredura}
\end{figure}

\subsection{Comparação de \emph{schedules}}

\begin{table}[H]
\centering
\caption{Tempo de execução (s) para diferentes \emph{schedules}, carga grande.}
\label{tab:schedules}
\small
\begin{tabular}{@{}lrrrr@{}}
\toprule
\emph{Schedule} & $T=2$ & $T=4$ & $T=8$ & $T=16$ \\
\midrule
\IfFileExists{plot/tab_schedules.tex}{
  \texttt{static} & 3{,}494086 & 2{,}002130 & 1{,}878103 & 1{,}895466 \\
\texttt{static, 1024} & 3{,}426964 & 1{,}956606 & 1{,}849952 & 1{,}883000 \\
\texttt{dynamic, 1024} & 3{,}673296 & 2{,}076709 & 1{,}930189 & 1{,}937384 \\
\texttt{guided} & 3{,}627628 & 2{,}065807 & 1{,}928899 & 2{,}041872 \\

}{
  \texttt{static} & 3,494086 & 2,002130 & 1,878103 & 1,895466 \\
  \texttt{static,1024} & 3,426964 & 1,956606 & 1,849952 & 1,883000 \\
  \texttt{dynamic,1024} & 3,673296 & 2,076709 & 1,930189 & 1,937384 \\
  \texttt{guided} & 3,627628 & 2,065807 & 1,928899 & 2,041872 \\
} \\
\bottomrule
\end{tabular}
\end{table}

As quatro políticas de escalonamento produzem tempos relativamente próximos entre si em todo o intervalo de $T$ testado, com uma diferença máxima de cerca de $8\%$ entre a mais rápida e a mais lenta em qualquer $T$ fixo (Tabela~\ref{tab:schedules}), o que indica baixo desbalanceamento de carga entre células: o corpo de trabalho por célula é majoritariamente uniforme (Seção~\ref{sec:omp}), sem concentrações de trabalho que justifiquem escalonamento dinâmico. Interessantes notar que quem "vence" nos quatro valores de $T$ testados não é o \texttt{static} padrão, e sim \texttt{static,1024}, o mesmo \texttt{schedule} estático, mas com \emph{chunk} fixo de 1024 elementos em vez do \emph{chunk} contíguo padrão ($\approx L \cdot C / T$ elementos), por margem de $2$ a $4\%$ sobre o \texttt{static} padrão, que fica em segundo lugar; um \emph{chunk} menor e fixo parece favorecer um pouco o aproveitamento de cache por \emph{thread} nesta carga, sem pagar o custo de sincronização por \emph{chunk} que \texttt{dynamic} paga a cada despacho. \texttt{dynamic,1024} e \texttt{guided} disputam as duas últimas posições sem um padrão único: \texttt{dynamic,1024} é a mais lenta em $T=2$ e $T=4$, \texttt{guided} em $T=16$, e as duas ficam muito próximas em $T=8$, o que, somado à margem estreita de \texttt{static,1024} sobre \texttt{static}, é consistente com uma diferença de desempenho pequena entre as quatro políticas, e não com uma política sistematicamente pior que as outras.

%=====================================================================
\section{Análise dos resultados}
\label{sec:analise}

O speedup bruto de $3{,}29$ a $3{,}54\times$ (Tabela~\ref{tab:speedup}) já é, ele mesmo, uma medida quase direta do paralelismo obtido. A Seção~\ref{sec:monothread} isolou, comparando as duas versões com $T=1$ em ambas (portanto sem nenhum trabalho paralelo real), um fator monothread de apenas $0{,}94$ a $0{,}97\times$, ou seja, próximo de 1, e nunca a favor da versão paralela: com uma única \emph{thread}, \texttt{fire\_omp} é tão rápido quanto \texttt{fire\_seq} na carga pequena e cerca de $6\%$ mais lento na carga média, resultado corroborado de forma independente por $S(T{=}1)=0{,}94$ na varredura de $T$ da carga grande (Seção~\ref{sec:varredura-t}). Isso é consistente com o estado atual do código, onde todas as otimizações sequenciais implementadas têm forma textualmente idêntica em \texttt{fire\_seq.c} e em \texttt{fire\_omp.c}; só a presença das diretivas \texttt{\#pragma omp} em torno delas muda entre os dois arquivos. Como os dois laços \emph{branchless} são curtos (oito iterações fixas), o resultado sugere que a diretiva \texttt{simd} explícita não é uma fonte de ganho adicional. A pequena desvantagem residual da versão paralela com $T=1$ é mais bem explicada pelo custo fixo de entrar em uma região \texttt{omp parallel} e sincronizar via \texttt{omp single} a cada passo (Seção~\ref{sec:omp}), mesmo sem paralelismo real, do que por qualquer otimização serial que ainda falte na versão sequencial.

Descontado esse fator quase neutro, o ganho atribuível apenas à paralelização é $S_\text{par} \approx 3{,}51$ a $3{,}64\times$ (Tabela~\ref{tab:speedup-corrigido}), muito próximo do speedup bruto (a correção nunca ultrapassa $7\%$). A eficiência correspondente depende de qual contagem de \emph{threads} se usa como referência: dividida pelos $T=4$ ou $T=8$ \emph{threads} nominais de cada entrada, a eficiência cai de $91\%$ (carga pequena, $T=4$) para $44\%$ (cargas média e grande, $T=8$), uma queda expressiva. Dividida pelos 4 núcleos físicos reais do processador (Tabela~\ref{tab:ambiente}), no entanto, a eficiência fica em $88$ a $91\%$ nas três cargas, próxima do linear. As duas leituras não são contraditórias, elas medem coisas diferentes. A primeira mostra que, ao usar 8 \emph{threads} nominais em um processador com apenas 4 núcleos físicos, a segunda metade das \emph{threads} (as lógicas, via \emph{hyperthreading}) contribui muito pouco além do que os 4 núcleos físicos já entregam sozinhos, o comportamento esperado de \emph{hyperthreading} em um laço com uso intenso de unidades aritméticas inteiras e pouca espera por memória. A segunda mostra que, contra o paralelismo físico real disponível, a estratégia de paralelização em si (região persistente, \texttt{omp for collapse(2)}, reduções) está próxima do ótimo. A leitura relevante para avaliar a qualidade da paralelização implementada é a segunda; a primeira é relevante para decidir quantas \emph{threads} de fato vale a pena usar neste processador, e a resposta, à luz da varredura de $T$ (Seção~\ref{sec:varredura-t}), é 4: o tempo praticamente estagna de $T=4$ para $T=8$ e $T=16$, e a eficiência bruta despenca no mesmo ponto.

%=====================================================================
\section{Conclusão}
\label{sec:conclusao}

As implementações sequencial e paralela produziram resultados numericamente idênticos em todos os cenários testados, três cargas de trabalho e cinco valores de $T$, confirmando que a paralelização por dados dentro de cada passo, sem quebra de dependência temporal entre passos, preserva exatamente a semântica do modelo sequencial (Seção~\ref{sec:validacao}). A estratégia de região paralela persistente, \texttt{omp for collapse(2)} e reduções inteiras eliminou a necessidade de sincronização explícita (\texttt{critical}/\texttt{atomic}) no laço principal, apoiando-se apenas nas barreiras implícitas do OpenMP para a correção da troca de \emph{buffers} e da condição de parada (Seção~\ref{sec:omp}).

A decomposição do speedup bruto (Seção~\ref{sec:analise}) mostrou que, com o código em seu estado atual, o fator monothread entre as duas versões está próximo de 1 ($0{,}94$ a $0{,}97\times$). Como resultado, o speedup bruto medido já é, essencialmente, o próprio ganho atribuível aos múltiplos núcleos, sem precisar de uma correção expressiva. Esse ganho é consistente com a arquitetura de 4 núcleos físicos do ambiente de testes (Tabela~\ref{tab:ambiente}): a eficiência paralela normalizada por esses 4 núcleos fica em $88$--$91\%$, ao passo que a mesma eficiência normalizada pelas 8 \emph{threads} nominais do processador cai para $44\%$ nas cargas média e grande, evidenciando o limite prático do \emph{hyperthreading} para esta carga de trabalho. A comparação entre políticas de escalonamento (\texttt{static}, \texttt{static,1024}, \texttt{dynamic,1024}, \texttt{guided}) mostrou diferenças moderadas entre si (no máximo cerca de $8\%$) em todo o intervalo de $T$ testado, com \texttt{static,1024} sendo a opção mais rápida em todos os casos testados, por margem estreita sobre o \texttt{static} padrão; ainda assim, a diferença pequena entre as quatro políticas indica que a carga de trabalho por célula é suficientemente uniforme para que o escalonamento estático seja uma escolha adequada, sem necessidade de balanceamento dinâmico.

\appendix

\section{Guia de compilação e execução}
\label{sec:apendice-execucao}

Esta seção resume o necessário para compilar e executar as duas versões a partir dos três arquivos entregues (\texttt{fire\_seq.c}, \texttt{fire\_omp.c} e \texttt{Makefile}), seguindo a forma de invocação exigida no enunciado (Seção 3).

\begin{lstlisting}[language=make,caption={Compilação.},basicstyle=\ttfamily\footnotesize]
make all      # compila as duas versoes: gera os binarios fire_seq e fire_omp
make clean    # remove os binarios gerados, se for preciso recompilar do zero
\end{lstlisting}

O alvo \texttt{all} chama o \texttt{gcc} com as \emph{flags} \texttt{-Wall -Wextra -O2 -fopenmp -std=c99} (Tabela~\ref{tab:ambiente}); não é necessário passar nenhuma flag manualmente.

\begin{lstlisting}[language=bash,caption={Execução, no formato exigido pelo enunciado.},basicstyle=\ttfamily\footnotesize]
./fire_seq entrada.txt
./fire_omp entrada.txt
\end{lstlisting}

Cada binário recebe apenas o caminho do arquivo de entrada como argumento e escreve as doze linhas de saída (Seção~\ref{sec:validacao}) em \texttt{stdout}. O número de \emph{threads} $T$ usado por \texttt{fire\_omp} é lido da própria primeira linha do arquivo de entrada (Seção~\ref{sec:comum}), não de uma variável de ambiente ou de um argumento adicional.

\end{document}